\begin{proposition}[Invariants of the compact torus fibration]\hypertarget{tag-02XO}{}\label{night:P4.5}\label{stmt:P4.5}
\StatementRecord{night:P4.5}{proved}{02XO}{\SeriesRef{night:P4.4}}

For the threefold $Y$ of Proposition~\ref{night:P4.2},
\[
 a(Y)=1,\qquad c_3[Y]=0,\qquad b_3(Y)=2,
\]
and the integral homology is free of ranks $(1,0,0,2,0,0,1)$.
\end{proposition}
\begin{proof}

\begin{enumerate}[label={$\langle1\rangle$\arabic*.},leftmargin=*]
\item The regular fibre has Euler number zero. Both components of each special
fibre are elliptic bundles over compact curves, and their intersection is
elliptic, so additivity gives zero for each special support. Stratifying
the base and applying multiplicativity to the smooth bundle gives
$e(Y)=0$. The Euler class of the real bundle underlying a complex
rank-three tangent bundle is its top Chern class, with the complex
orientation. Hence $c_3[Y]=0$. Duality and Proposition~\ref{night:P4.4} give $b_4=b_5=0$ and
$e(Y)=2-b_3$, so $b_3=2$. Integral duality and the universal-coefficient
theorem identify the torsion of $H_3$ with that of $H_2$, which is zero.
All other groups have already been determined.

\item We also have $H^2(Y;\Z)=0$. A horizontal effective divisor would restrict
to a nonzero effective divisor on the Kähler torus $A$ with zero first
Chern class. Its pairing with a Kähler class would be positive by
integration over the divisor and zero by the Chern-class computation.
Thus every divisor on $Y$ is vertical. The polar divisor of a meromorphic
function is consequently vertical; on a general fibre the function is
holomorphic and therefore constant. Local holomorphic sections give a
descended meromorphic function on the regular part of the base. Near
each exceptional base value a sufficiently high power of the base
coordinate removes the vertical poles. Connected fibres then give
holomorphic descent there, and dividing by that power gives meromorphic
descent. Therefore $\C(Y)=f^*\C(\Pp^1)$ and $a(Y)=1$.

\end{enumerate}
\end{proof}